\section{Level loss and a convex singular-value weight}
\label{rec:sec-grid}
\label{sec:convex-grid}

The final three sections give alternative recursions across a balanced
coordinate split. Their final indexed power bounds are equivalent up
to the choice of an absolute exponent. The estimates preserved during
the induction carry different information: the first records the
level lost on restriction to the lines, the second keeps the ranks and
norms of local Fourier bands, and the third supplies a pointwise
smoothing bound that remains valid after conditioning a completed
grid on compatibility.

Put $a=2^{\lfloor d/2\rfloor}$, $b=2^{\lceil d/2\rceil}$ and
$n=ab$. The positions form an $a$ by $b$ grid, and $T_n=YX$, where
$X$ is the tensor product of the row sweeps and $Y$ the tensor
product of the column sweeps. At this split the analytic difficulty
is the overlap between the row and column Fourier subspaces.
We estimate that overlap before applying the decay available in
the smaller sweeps.

\subsection{The weight and the level deficit}

For $\lambda=(n-k,\beta)$ and $1\le j\le D_\lambda$, put
$M=k\log n$, $D=\log(D_\lambda j)$, and
\begin{equation}\label{rec:perspective-weight}
 \mathcal W(M,D)=M\max\{10^{-8}D/M,(D/M)^3\},\qquad
 \mathcal W(0,0)=0.
\end{equation}
For $k>0$, the dimension bound $D_\lambda\le n^k$ gives
$0\le D/M\le2$; the trivial type uses the convention at $(0,0)$.
The linear term prevents the weight from becoming too small at
low dimension. The cubic term compensates for loss of level on
restriction: at fixed $D$, the expression $D^3/M^2$ increases when
the available level mass $M$ decreases. To verify the property needed in the recursion,
write $f(u)=\max\{10^{-8}u,u^3\}$. It is convex on $[0,\infty)$
and $f(0)=0$. If $M>0$ and $M'=\sum_iM_i\le M$, Jensen's inequality, with
an additional term of weight $M-M'$ and argument zero, gives
\begin{equation}\label{grid:perspective-jensen}
 \sum_i\mathcal W(M_i,D_i)
 =\sum_iM_i f(D_i/M_i)
 \ge M f\left(\frac{\sum_iD_i}{M}\right).
\end{equation}
Terms with $M_i=0$ have $D_i=0$ and contribute zero. On the range
$0\le D/M\le2$ we also have
$10^{-8}D\le\mathcal W(M,D)\le4D$. We will preserve the full
weight through the proof, since this comparison alone does not
explain how it closes under restriction.

\begin{theorem}\label{rec:perspective-singular}
There is an absolute $c>0$ such that, for every dyadic $n\ge2$,
every nontrivial partition $\lambda\vdash n$, and
$1\le j\le D_\lambda$,
\[
 s_j(T_n|_{V_\lambda})\le
       \exp\{-c\mathcal W(k\log n,\log(D_\lambda j))\}.
\]
In particular all singular values are bounded by a fixed negative power
of $D_\lambda j$.
\end{theorem}

Here and below representations with more than $n/2$ rows are killed by
the first matching average.  We therefore omit them when using dimension
lower bounds.

For each direction, sum $q-\rho_1$ over its line types
$\rho\vdash q$, and call these total levels $k_A,k_B$.
The second-overlap estimate below incurs the explicit loss
\[
 \frac{\log(n/k)}2\left(k-\frac{k_A+k_B}{2}\right).
\]
The cubic perspective compensates for this term. Large-dimensional
parent blocks instead use a fourth-overlap estimate, and the smallest
levels are controlled directly on a larger tuple module. We establish
these inputs before closing the induction. The interpolation at the
end of the argument will allow an arbitrarily small positive exponent
from the finite base cases.

\subsection{Fourth and second overlap bounds}

\begin{lemma}\label{rec:biased-grid-fourth}
There is an absolute constant $C>0$ such that the following holds
for every $n=2^d$ with $d\ge1$, on the balanced grid
$a=2^{\lfloor d/2\rfloor}$, $b=2^{\lceil d/2\rceil}$.
For each row or column line $i$, choose an irreducible type $\rho_i$
of its line symmetric group. Let $V_{\rho_i}$ be one carrier, put
$D_{\rho_i}=\dim V_{\rho_i}$, and choose an orthogonal projection
$E_i$ on $V_{\rho_i}$, extended by zero on the other Fourier types.
Let $R$ and $S$ be the tensor products of these projections over
the rows and columns, respectively, viewed as group-algebra operators
for $S_n$. If any $E_i$ is zero, then $RS=0$. Otherwise put
\[
 r_i=\operatorname{rank}_{V_{\rho_i}}E_i,\qquad
 D_c=\sum_i\log(D_{\rho_i}r_i),
\]
where the sum includes both directions. Thus $r_i$ is a positive
carrier rank, and the corresponding line group-algebra projection
has rank $D_{\rho_i}r_i$ in the regular representation.
For every $\lambda\vdash n$, write $R_\lambda,S_\lambda$ for the
actions on one copy of $V_\lambda$. Then, with unnormalized trace,
\[
 D_\lambda\Tr_{V_\lambda}
   (R_\lambda S_\lambda R_\lambda S_\lambda)
 \le\exp\{D_c+C D_c/\log n+C n^{1-1/40}\}.
\]
The factor $D_\lambda$ is the multiplicity in the regular
representation. No restriction is imposed on the parent partition
$\lambda$.
\end{lemma}
\begin{proof}
Let $K$ and $L$ be the row and column permutation subgroups.
A line projection of type $\rho$ and carrier rank $r$ has regular
rank $D_\rho r$. Its coefficient density $f$ therefore satisfies
$\E|f|^2=D_\rho r$, and $|f|^2/(D_\rho r)$ is a probability
density bounded by $D_\rho r$, by
Lemma~\ref{lem:coefficient-compatibility}. Applying that lemma to
the tensor products gives
\[
 D_\lambda\Tr_{V_\lambda}
       (R_\lambda S_\lambda R_\lambda S_\lambda)
 \le\Tr_{\rm reg}(RSRS)\le
 e^{D_c}\frac{n!}{|K||L|}\Pr(\mathcal S),
\]
where row and column permutations are independent, their line laws
are these squared-coefficient densities, and $\mathcal S$ is their
compatibility event. The lemma includes the unique-completion and
inversion changes that identify this event with the four-factor trace.

Fix the column permutations.  Write $j_i(t)$ for the original column
in row $i$ sent by the row permutation to column $t$, and $k_t(i)$
for the subsequent output row under the column permutation.  Success
requires the pairs $(k_t(i),j_i(t))$, over all cells $(i,t)$, to be
distinct.  Put $r_{ik}^{\circ}=|\{t:k_t(i)=k\}|$.  Call a cell bad
if $r_{i,k_t(i)}^{\circ}>b^{1/8}$, and let $H$ count the bad cells.
This designation depends only on the fixed column permutations.

Let $\rho_i$ be the original law of row $i$, with density at most
$K_i$ relative to its uniform law $U_i$.  Put
$p=\Pr_{\otimes_i\rho_i}(\mathcal S)$ and assume $p>0$; the
zero-probability case needs no estimate.  Let
$Q=(\otimes_i\rho_i)(\,\cdot\mid\mathcal S)$, and write $Q_i$
for its row marginals.  Define
\[
 I=\KL(Q\Vert\otimes_iQ_i),\qquad
 D=\sum_i\KL(Q_i\Vert\rho_i),\qquad
 L_i=\KL(Q_i\Vert U_i).
\]
The relative-entropy chain rule and the density caps give the exact
identity and inequality
\begin{equation}\label{rec:biased-success-KL}
 -\log p=I+D,\qquad
 \sum_iL_i\le D+\sum_i\log K_i.
\end{equation}
Indeed $\KL(Q\Vert\otimes_i\rho_i)=-\log p$, and
$L_i=\KL(Q_i\Vert\rho_i)+
\E_{Q_i}\log(d\rho_i/dU_i)$.  These formulas retain the extra
relative entropy created by conditioning on success.

Reveal rows in independent continuous uniform priority order, and
within each row use an independent uniform order of its cells.  At
a cell of row $i$, let $P(j)$ be the prediction under $Q_i$ given
only the already exposed cells of that row.  A symbol is called
heavy at this revelation when $P(j)>b^{-1/2}$.  Heavy symbols are
different from the bad cells defined above.  Except at the last
$\lceil b^{4/5}\rceil$ revelations of a row, the uniform permutation
prediction assigns each available symbol probability at most
$b^{-4/5}$.  On a heavy symbol, the nonnegative divergence term
$u\log(u/v)-u+v$ is at least $c u\log b$.  Summing the prediction
divergences by the chain rule therefore gives
\begin{equation}\label{rec:biased-heavy-count}
 \E_Q\#\{\text{heavy outcomes}\}
 \le C ab^{4/5}+\frac{C}{\log b}\sum_iL_i.
\end{equation}

Consider a good cell $(i,t)$ and write $y=1-\text{priority}(i)$.
In the chain-rule expression for $I$, the reference prediction is
$P$, whereas the actual prediction conditions on the full past.
Only symbols not already used at output row $k_t(i)$ are allowed.
Use the unnormalized reference
$P_*(j)=yP(j)$ on heavy symbols and $P_*(j)=P(j)$ on the others.
If $A$ is the set of allowed symbols and $q$ is the actual conditional
prediction, normalization of $P_*\mathbf1_A$ gives
\[
 \KL(q\Vert P)\ge
 -\log\sum_{j\in A}P_*(j)+q(\text{heavy})\log y.
\]
Conditional on the complete successful assignment, the within-row
order, and $y$, every symbol at this output row belongs to exactly
one packet.  If that packet comes from a different input row, its
row has not yet been exposed with probability $y$.  There are at
most $b^{1/8}$ exceptions coming from the current row.  Their light
symbols contribute at most $b^{1/8}b^{-1/2}$; their heavy symbols
are covered by the factor $y$.  Thus the conditional expected
allowed mass is at most $y+b^{-3/8}$.  Jensen's inequality and
integration over $y\in[0,1]$ make the first term at least
$1-O(b^{-1/4})$.  The expected cost of the second term is the
heavy-outcome probability, since the assignment and within-row
order are independent of the row priority and
$\int_0^1\log y\,dy=-1$.

For bad cells retain the nonnegative raw conditional KL contribution
and charge no unit saving.  Summing good cells and using
\eqref{rec:biased-heavy-count} yields
$I\ge n-H-Cn^{1-1/40}-a_b\sum_iL_i$, where $a_b=C/\log b$.
Consequently \eqref{rec:biased-success-KL} gives
\begin{equation}\label{rec:biased-entropy-absorption}
 -\log p\ge n-H-Cn^{1-1/40}
       -a_b\sum_i\log K_i+(1-a_b)D.
\end{equation}
For sufficiently large $b$, $a_b<1$.  Discarding the nonnegative
last term proves the conditional success estimate
\begin{equation}\label{rec:biased-grid-success}
 \log p\le-n+H+Cn^{1-1/40}
                       +\frac{C}{\log b}\sum_i\log K_i.
\end{equation}
In particular the main saving remains exactly $-n$.

It remains to average the factor $e^H$ over the biased columns.
We give the required bundle count. Put $B=\lfloor b^{1/8}\rfloor+1$.
A bundle chooses an input row $i$, a destination row $k$, and $B$
columns in which $k_t(i)=k$. A family of bundles is disjoint if
its prescribed cells are disjoint. Under independent uniform
column permutations, a fixed family of $s$ bundles has probability
at most $(e/a)^{Bs}$: a column with $u$ consistent prescribed
images has probability $1/(a)_u\le(e/a)^u$.
For each repeated row value with count $v\ge B$, pack
$\lfloor v/B\rfloor$ disjoint bundles. Thus every outcome has a
satisfied disjoint family with $2Bs\ge H$.
For $c=1/64$, summing over all possible such families, including
the empty one, proves
\[
 \mathbb E_U b^{cH}
 \le\sum_{s\ge0}
 \left\{a^2\binom bB(e/a)^B b^{2cB}\right\}^{s}
 \le2
\]
for sufficiently large $b$. Indeed the expression in braces is
at most $a^2(2e^2b^{2c}/B)^B$, which tends to zero.
The joint biased column law has a density $\rho$ bounded by
$e^{x_C}$, where $x_C$ is the column part of $D_c$, and has mean
one. H\"older with $r=c\log b>1$ yields
\[
 \mathbb E_U\rho e^H
 \le(\mathbb E_U\rho^{r/(r-1)})^{(r-1)/r}
       (\mathbb E_U e^{rH})^{1/r}
 \le\exp\{(x_C+\log2)/(c\log b)\}.
\]
Together with \eqref{rec:biased-grid-success}, this supplies a
compatibility probability at most
$\exp\{-n+Cn^{1-1/40}+CD_c/\log n\}$.
The exact $-n$ cancels the factorial factor, since
$\log(n!/(|K||L|))=n+O((a+b)\log n)$.
Multiplication by $e^{D_c}$ from the squared coefficient masses
proves the claimed bound for all sufficiently large $b$.
For bounded $n\ge2$, the same constant can be enlarged uniformly
over all local projections. Indeed $R_\lambda,S_\lambda$ are
orthogonal projections, so
\[
 D_\lambda\Tr(R_\lambda S_\lambda R_\lambda S_\lambda)
 \le D_\lambda^2\le n!.
\]
Since $D_c\ge0$, a fixed multiple of $n^{39/40}$ absorbs
$\log(n!)$ over the finitely many remaining sizes. This includes
the grid $(a,b)=(1,2)$; a one-site line has only the trivial
carrier. No estimate with denominator $\log1$ is needed.
\end{proof}

The middle range needs a different estimate.  If
$n^{3/5}<k<n^{1-1/250}$, set $b_0=\log(n/k)$, and let $k_A,k_B$ be
the total row and column levels of the product projections. We may
assume both types occur in the restriction of $V_\lambda$, since
otherwise the overlap is zero. The Littlewood--Richardson rule then
gives $k_A,k_B\le k$. If $\Tr(RS)=0$, the exponentiated inequality
below is automatic. Otherwise assume $\Tr(RS)>0$, so that its
logarithm is defined. We claim
\begin{equation}\label{rec:tuple-grid-second}
 \log\{D_\lambda\Tr(RS)\}
 \le D_c+\frac{b_0}{2}(k-k_A/2-k_B/2)
                  +Ck(\log n)^{3/4}.
\end{equation}
To prove this, take $l=\lceil\sqrt{nk}\rceil$ labelled marks and put
$z=l/n$. By branching, the multiplicity of $V_\lambda$ in this
tuple module is $\binom lkD_\beta$; the first-row strip and the lower
diagram occupy disjoint columns. Moreover
\[
 \binom lkD_\beta\ge D_\lambda z^k e^{-Ck},
\]
because $D_\lambda\le\binom nkD_\beta$ and
$\binom lk/\binom nk\ge(l/n)^k e^{-Ck}$ in this range.

Choose now the $l$ marked sites independently and uniformly, allowing
repetitions, and let $U$ be the event that they are distinct. Let
$u_i$ and $v_j$ be their row and column counts. In a trace term for
a row move followed by a column move, every mark must be fixed
individually. Since the second move cannot change its column, the
first move must fix it; the second then fixes it as well. Thus only
the pointwise stabilizers of the marked sites contribute.

For a row projection of type $\rho$ and carrier rank $r$, write
$P_\rho$ for that carrier projection. If $A$ is the set of marked
sites in the row, put
\[
 h_\rho(A)=r^{-1}\operatorname{Tr}
       (P_\rho\Pi_{S_{[b]\setminus A}}),
\]
where $\Pi$ denotes averaging over the displayed pointwise
stabilizer. This lies in $[0,1]$. The sum of the projection's
group-algebra coefficients over that stabilizer is
$D_\rho r\,h_\rho(A)/(b)_{|A|}$. Multiplying these formulas in
both directions gives the exact tuple trace identity
\[
 \operatorname{Tr}_{[l]}(RS)
 =e^{D_c}\mathbb E\{\mathbf1_Uw_Aw_Bh_Ah_B\},\qquad
 w_A=\prod_i\frac{b^{u_i}}{(b)_{u_i}},\quad
 w_B=\prod_j\frac{a^{v_j}}{(a)_{v_j}}.
\]
Here $\operatorname{Tr}_{[l]}$ is the trace on ordered injective
$l$-tuples, and $h_A,h_B$ are the products of the corresponding
local overlaps. Outside the event that all counts are admissible
set the weights to zero; only $U$ contributes to the identity.

We need two bounds for this expectation. First, if a row type has
level $h=b-\rho_1$, then averaging its stabilizer projection over
uniform marked sets of size $u$ gives a scalar projection average.
The scalar is the invariant dimension divided by $D_\rho$.
Branching bounds the numerator by
$\binom uhD_{\bar\rho}$, while
\[
 D_\rho\ge2^{-h}\binom bhD_{\bar\rho}.
\]
For the latter inequality, construct the lower diagram one box at a
time; each added box increases a top-row hook by a factor at most
two. Averaging over the multinomial row counts therefore gives
\[
 \mathbb E(\mathbf1_Uh_A)
 \le\prod_i\frac{2^{h_i}}{(b)_{h_i}}
        \mathbb E\prod_i(u_i)_{h_i}
 \le (2e)^{k_A}z^{k_A}.
\]
We used $\sum_i h_i=k_A$,
$\mathbb E\prod_i(u_i)_{h_i}=(l)_{k_A}a^{-k_A}$, and
$(b)_h\ge(b/e)^h$. Conditional on row counts, dropping the
probability of distinct marked columns can only increase this
bound. The column argument gives
$\mathbb E(\mathbf1_Uh_B)\le(2e)^{k_B}z^{k_B}$.

Second, with $r=(\log n)^{1/4}$,
\begin{equation}\label{grid:occupancy-weight}
 \log\|w_Aw_B\|_r\le C\{l^2/n+(a+b)\log n\}.
\end{equation}
Here are details of the occupancy bound. For the row counts, first
use independent Poisson variables of mean $\lambda=bz$ and
condition their sum to be $l$. This produces the multinomial law,
and the conditioning costs only $O(\log n)$ in the logarithm of
the moment. For $0\le u\le b$, Stirling's bounds give
\[
 \log\frac{b^u}{(b)_u}
 \le b\phi(u/b)+C\log(b+1),\qquad
 \phi(s)=s+(1-s)\log(1-s).
\]
Use $\phi(s)\le Cs^2$ for $s\le1/2$ and $\phi(s)\le s$
throughout $[0,1]$. The Poisson rate function is
$I_\lambda(u)=u\log(u/\lambda)-u+\lambda$. For every
$1\le s_0\le2r$, and all sufficiently large $n$, these inequalities
imply
\[
 s_0b\phi(u/b)
 \le C s_0\lambda^2/b+\tfrac12I_\lambda(u).
\]
Indeed $u\le C\lambda$ is paid by the first term. For larger
$u\le b/2$, compare $s_0u^2/b$ with
$u\log(u/\lambda)$; their ratio is small uniformly because
$s_0z=o(1)$ and $s_0/\log(1/z)=o(1)$. For $u>b/2$, the latter
condition compares $s_0u$ with the same rate function. Summing
the Poisson probabilities times the weight to power $s_0$ now
costs at most
$C s_0\lambda^2/b+C s_0\log n$ in its logarithm, since there
are at most $b+1$ admissible counts. Product over rows and
conditioning the total give
$\log\|w_A\|_{2r}\le C(l^2/n+a\log n)$.
The column bound is analogous, and H\"older proves
\eqref{grid:occupancy-weight}; no independence between row and
column counts is required.

Finally apply H\"older with exponents $r,2r/(r-1),2r/(r-1)$.
Since $0\le h_A,h_B\le1$, the logarithm of the tuple trace is at
most
\[
 D_c+C\{k+(a+b)\log n\}
 +\frac{1-r^{-1}}2(k_A+k_B)\log z+C(k_A+k_B).
\]
All irreducible contributions to $\operatorname{Tr}_{[l]}(RS)$
are nonnegative. Divide by the multiplicity lower bound above.
As $k_A,k_B\le k$, $-\log z=b_0/2+O(1)$, and
$k>n^{3/5}$, the resulting error is bounded by
$Ck(\log n)^{3/4}$. This is \eqref{rec:tuple-grid-second}.

\subsection{Small levels and power-product interpolation}

The small-level input is an explicit contraction, rather than a
finite-dimensional base case:
\begin{equation}\label{rec:perspective-small-level}
 k\le n^{3/5}\quad\Longrightarrow\quad
 \|T_n|_{V_\lambda}\|_{\op}\le e^{-c k\log n}.
\end{equation}
Use $l=\lceil k n^{1/50}\rceil$ and the permutation module on ordered
injective $l$-tuples.  Denote its sweep operator by $T_n^{[l]}$.
For a fixed realization $\pi$ and an $l$-set $I$ of initial sites,
the endpoint of each tracked card determines its entire route through
the sweep: each bit takes its final value at its unique update.
An independent realization $\pi'$ therefore agrees on all endpoints
with probability $n^{-l}2^{e(I)}$, where $e(I)$ counts switches
visited by two tracked cards.  The requirements are consistent
because $\pi$ realizes them.  Summing the squared transition
probabilities over ordered initial and final tuples gives
\begin{equation}\label{rec:perspective-tuple-HS}
 \|T_n^{[l]}\|_{\HS}^2
 =\frac{(n)_l}{n^l}\E_{\pi,I}2^{e(I)}
 \le\E_{\pi,I}2^{e(I)},
\end{equation}
where $I$ is uniform among $l$-subsets.  This is the full tuple
Hilbert--Schmidt norm; no centering is needed here.

For fixed wiring, the co-visit graph has maximum degree at most $d$.
Its induced graph on any $s$ sites has at most $s\log_2s/2$ edges.
Indeed the recursive split into two subcubes adds at most
$\min(s_0,s_1)$ edges, and
$s_0\log_2s_0+s_1\log_2s_1+2\min(s_0,s_1)
\le(s_0+s_1)\log_2(s_0+s_1)$.
There are at most $nd^{2s}$ connected vertex sets of size $s$, by
encoding a root and a walk along a spanning tree.

Expand an upper bound using unordered collections of disjoint
connected subsets, each of size between $2$ and $l$, weighted by
the product of $s^{s/2}$.  For a given $I$, its actual nontrivial
components occur in this sum and their product bounds $2^{e(I)}$.
A collection on $u$ sites has inclusion probability
$(l)_u/(n)_u\le(l/n)^u$.  After taking expectation, dropping
disjointness and then using $1+x\le e^x$ gives the finite expansion
\begin{equation}\label{rec:perspective-finite-encounters}
 \log\E 2^{e(I)}
 \le\sum_{s=2}^{l}nd^{2s}(l/n)^s s^{s/2}
 =O(d^4l^2/n).
\end{equation}
Consecutive summands have ratio at most $C d^2(l/n)\sqrt{l+1}$.  Since
$l\le2n^{31/50}$ in the stated range, this is
$O(d^2n^{-7/100})=o(1)$.  The finite sum is thus bounded by a
constant times its $s=2$ term.  No extension to arbitrarily large
$s$ is made.

The tuple multiplicity of $V_\lambda$ is
$f^{\lambda/(n-l)}$ by branching.  For all sufficiently large $n$,
$n-l\ge k\ge\lambda_2$, so the remaining first-row strip of length
$l-k$ and the lower diagram $\beta$ occupy disjoint columns.
Their tableau entries can be interleaved freely, giving
\begin{equation}\label{rec:perspective-tuple-multiplicity}
 f^{\lambda/(n-l)}=\binom lk D_\beta.
\end{equation}
The tuple action is the orthogonal sum of its irreducible blocks,
with these multiplicities.  Positivity of their Hilbert--Schmidt
squares, \eqref{rec:perspective-tuple-HS}, and
\eqref{rec:perspective-finite-encounters} imply
\[
 \|T_n|_{V_\lambda}\|_{\op}^2
 \le\frac{\exp(Cd^4l^2/n)}{\binom lk D_\beta}.
\]
Now $\binom lk\ge(l/k)^k\ge n^{k/50}$, while
$d^4l^2/n=o(k\log n)$ uniformly for $k\le n^{3/5}$.
Taking square roots proves \eqref{rec:perspective-small-level},
for example with any fixed $c<1/100$ after increasing the absolute
starting dimension.

We first record the interpolation step that allows the exponent in the
induction hypothesis to be arbitrarily small.  Its conclusion concerns
products of singular values, which is why the last step of the proof
below treats all indices up to the requested index.

\begin{lemma}\label{grid:perspective-power-interpolation}
If $A,C$ are positive semidefinite operators on a finite-dimensional
Hilbert space and $v\ge1$, then, for every admissible $j$,
\begin{equation}\label{grid:perspective-power-products}
 \prod_{i=1}^j s_i(CA)
 \le \left(\prod_{i=1}^j s_i(C^vA^v)\right)^{1/v}.
\end{equation}
\end{lemma}
\begin{proof}
This is Lemma~\ref{lem:power-product-comparison}, including its
continuity argument at zero eigenvalues.
\end{proof}

\subsection{Closing the convex-weight induction}

\begin{proof}[Proof of Theorem~\ref{rec:perspective-singular}]
Write $\eta=10^{-8}$ and $u_0=3/1000$.  Fix a constant
$c_s\in(0,1/100)$ for which
\eqref{rec:perspective-small-level} holds above an absolute cutoff,
and put
\[
 c_0=\min\{1/1000,c_s/16\}.
\]
In particular $12c_0<1/4$.  We will specify an absolute starting
exponent $d_0$ after estimating the errors; its choice will depend
on $c_0$ and the absolute constants in the two overlap estimates,
but not on the exponent supplied by the finite base cases.

Induct on $d=\log_2n$.  Denote the available exponents for the two
child sizes by $c_{\lfloor d/2\rfloor}$ and
$c_{\lceil d/2\rceil}$, and set
\[
 c=\min\{c_{\lfloor d/2\rfloor},c_{\lceil d/2\rceil}\},
 \qquad v=c_0/c.
\]
All these exponents will lie in $(0,c_0]$, so $v\ge1$.  Fix a
nontrivial $\lambda=(n-k,\beta)$ that is not annihilated and put
$M=k\log n$.  Since $\mathcal W(M,D)\le4D\le8M$, the range
$k\le n^{3/5}$ already follows from
\eqref{rec:perspective-small-level}, with any exponent at most
$c_0$.  We henceforth assume $k>n^{3/5}$.

Here is the band decomposition used at a split.  For a child size
$m\in\{a,b\}$ and a type $\rho\vdash m$, let
$h=m-\rho_1$, $q=D_\rho$, and $M_\rho=h\log m$.
Choose an orthonormal eigenbasis, in decreasing eigenvalue order, of
the appropriate absolute value $|T_m|$ or $|T_m^*|$ of the child
sweep.  For each integer
$0\le t\le\lfloor\log_2q\rfloor$, the $t$th band is the span of
the eigenvectors with indices
\[
 I_t=\{2^t,\ldots,\min(2^{t+1}-1,q)\}.
\]
Its starting index $r=2^t$ bounds its carrier rank from above.
Associate to it the number $D_\rho(r)=\log(qr)$.  The induction
hypothesis bounds the operator norm on this band by
\[
 \exp\{-c\mathcal W(M_\rho,D_\rho(r))\}.
\]
For the trivial type there is one band, with $M_\rho=D_\rho(1)=0$.
An annihilated type contributes the zero operator and may be omitted.
The bound $q\le m^h$ shows that every band satisfies
$0\le D_\rho(r)\le2M_\rho$.

A row profile specifies a type and a band on each of the $a$ rows;
a column profile does the same on each of the $b$ columns.  We keep
only product types occurring in the restriction of $V_\lambda$.
If their total levels are $k_A$ and $k_B$, respectively, the
Littlewood--Richardson rule gives $k_A,k_B\le k$.  Indeed a
constituent of an induced product of the local types has first row
at most the sum of their first rows.  For a pair of profiles write
\begin{equation}\label{grid:perspective-profile-parameters}
 \begin{split}
 D_c&=\sum_i D_i,\qquad
 F_c=\sum_i\mathcal W(M_i,D_i),\\
 M'&=\sum_iM_i=k_A\log b+k_B\log a\le k\log n=M.
 \end{split}
\end{equation}
The sums include both directions.  Thus $D_c\le2M'\le2M$, and
\eqref{grid:perspective-jensen} gives
$F_c\ge\mathcal W(M,D_c)$.

We will need a count of these profiles before using any decay
exponent.  There are at most $(k+1)^a$ row-level lists.
For any such list $(h_i)$ with $\sum h_i\le k$, the number of
row types is at most
\[
 \prod_{i=1}^a p(h_i)
 \le\exp\left(3\sum_{i=1}^a\sqrt{h_i}\right)
 \le\exp(3\sqrt{ak}).
\]
Here $p(0)=1$ and $p(h)\le e^{3\sqrt h}$ for $h\ge1$.
On each line the number of bands is at most
$1+\log_2(m!)\le n^2$ for $n\ge4$.  Since
$\log(k+1)\le2\log n$, the number of row profiles is at most
$\exp(4a\log n+3\sqrt{ak})$.  The same count for columns shows
that the number $L$ of pairs of nonzero profiles satisfies
\begin{equation}\label{grid:perspective-profile-count}
 \log L\le B,
 \qquad B=4(a+b)\log n+3(\sqrt a+\sqrt b)\sqrt k.
\end{equation}
If there are no nonzero pairs, the required operator is zero and
the conclusion is immediate.  Hence assume $L\ge1$.

Take polar decompositions in the row and column group algebras,
writing $X=A U_X$ and $Y=U_Y C$, where
$A=|X^*|$ and $C=|Y|$.  The partial isometries can be extended to
unitaries in each local carrier space.  Thus $T_n=U_YCAU_X$ has
the singular values of $CA$.  Put $Z=C^vA^v$.  If $R$ and $S$
are the projections of a row and a column profile, respectively,
then
\[
 Z=\sum_{(R,S)} C^v S R A^v.
\]
The projections commute with the positive factors in their own
direction.  The local band estimates therefore give, for every
index $r\ge1$,
\begin{equation}\label{grid:perspective-profile-singular}
 s_r(C^vSRA^v)\le e^{-c_0F_c}s_r(SR).
\end{equation}
This also follows directly by factoring out the two band norms
and using $s_r(B_1HB_2)\le\|B_1\|_{\op}\|B_2\|_{\op}s_r(H)$.
If a band's actual rank is less than its starting index, the
overlap estimates are still applicable: their actual parameter
$\sum_i\log(D_{\rho_i}\operatorname{rank}R_i)$ is at most
$D_c$, and each displayed upper bound increases with that
parameter.

For $1\le j\le D_\lambda$ set
\[
 D=\log(D_\lambda j),\qquad r_j=\lceil j/L\rceil.
\]
The singular-value sum inequality, or rank-$(r_j-1)$
approximations to all $L$ summands, gives
\begin{equation}\label{grid:perspective-rank-allocation}
 s_j(Z)\le L\max_{(R,S)}s_{r_j}(C^vSRA^v),
 \qquad \log r_j\ge\log j-B.
\end{equation}
In fact the sum of those approximations has rank at most
$L(r_j-1)<j$, which verifies the index in this inequality even
when $j<L$.

First suppose $\log D_\lambda\ge n^{99/100}$.  Since
$\Tr(RSRS)=\|SR\|_{S^4}^4$, Lemma~\ref{rec:biased-grid-fourth}
implies
\[
 s_r(SR)\le
 \exp\left[-\frac14\left\{
 \log(D_\lambda r)-D_c-\frac{C_4D_c}{\log n}
                         -C_4n^{39/40}\right\}_+\right]
\]
for an absolute $C_4$.  Combining this with
\eqref{grid:perspective-profile-singular} and
\eqref{grid:perspective-rank-allocation} yields the indexed bound
\begin{equation}\label{grid:perspective-dense-indexed}
 \log s_j(Z)\le B-\min_{(R,S)}\left[
 c_0F_c+\frac14\left\{
 D-D_c-\frac{C_4D_c}{\log n}-C_4n^{39/40}-B
 \right\}_+\right].
\end{equation}
As usual a zero singular value satisfies every such upper bound.

The function $f(u)=\max(\eta u,u^3)$ is increasing and is
$12$-Lipschitz on $[0,2]$.  If $D_c\ge D$, Jensen's inequality
already gives $F_c\ge\mathcal W(M,D)$.  Otherwise put
$x=D-D_c>0$ and
\[
 E_4=\frac{C_4D}{\log n}+C_4n^{39/40}+B.
\]
Jensen and the Lipschitz bound give
$F_c\ge\mathcal W(M,D)-12x$.  Since $12c_0\le1/4$,
\[
 -12c_0x+\tfrac14(x-E_4)_+\ge-12c_0E_4.
\]
Thus both cases of $D_c$ in
\eqref{grid:perspective-dense-indexed} imply
\begin{equation}\label{grid:perspective-dense-error}
 \log s_j(Z)\le-c_0\mathcal W(M,D)+B+12c_0E_4.
\end{equation}
To quantify this error, use $k\le n$, $a+b\le3\sqrt n$ and
$\sqrt a+\sqrt b\le3n^{1/4}$.  Above an absolute cutoff,
$B\le C_B n^{3/4}$ for an absolute $C_B$.  Since
$\mathcal W(M,D)\ge\eta D$ and $D\ge n^{99/100}$, there is
an absolute $K_4$, depending only on the already fixed $c_0,\eta$
and $C_4$, such that
\begin{equation}\label{grid:perspective-dense-relative-error}
 \frac{B+12c_0E_4}{c_0\mathcal W(M,D)}
 \le K_4\left((\log n)^{-1}+n^{-3/200}+n^{-6/25}\right).
\end{equation}
In particular none of these constants depends on $c$.

We next justify exactly the dimension range in the other case.
Suppose $\log D_\lambda<n^{99/100}$.  The surviving shape has
height at most $n/2$.  It must have first row longer than $n/2$
once $n$ is sufficiently large.  Indeed, otherwise both its width
and height are at most $n/2$.  Transpose if necessary so that the
first row has length $q=\max(\lambda_1,\lambda'_1)$.  If
$q\le n/8$, every hook has length at most $2q\le n/4$, and the
hook formula gives $D_\lambda\ge(4/e)^n$.  If $q>n/8$, retain
that first row and a subdiagram of $t=\lfloor q/4\rfloor$ tail
boxes.  Such a subdiagram exists because $n-q\ge n/2\ge t$.
Its tableaux extend to the full shape.  Applying
\eqref{eq:near-row-hooks} to the retained diagram gives
\[
 D_\lambda\ge\binom{q+t}{t}
       \exp\left(-\frac{t}{q-t+1}\right)
 \ge4^t e^{-1/3}.
\]
For large $n$, $t\ge q/8>n/64$.  Both alternatives give
$\log D_\lambda\ge c_1n$ for an absolute $c_1>0$, contradicting
the assumed upper bound.  Hence $k<n/2$.

If nevertheless $k\ge n^{249/250}$, retain the first row of
length $q=n-k>n/2$ and $t=\lfloor n^{249/250}\rfloor$ tail
boxes.  For sufficiently large $n$ we have $t\le q/2$, so the
same hook estimate implies
\[
 \log D_\lambda\ge t\log(q/t)-1
 \ge\frac1{2000}n^{249/250}\log n
 >n^{99/100}.
\]
For the middle inequality, use $t\ge n^{249/250}/2$ and
$\log(q/t)\ge(\log n)/250-\log2\ge(\log n)/500$,
and increase the cutoff to absorb the subtracted $1$.
We have proved $k<n^{249/250}$.  Moreover
\eqref{eq:near-row-hooks}, now applied to $\lambda$ itself,
gives
\begin{equation}\label{grid:perspective-middle-u}
 \log D_\lambda\ge k\log(n/k)-1,
 \qquad
 u:=\frac{D}{M}\ge\frac{\log(n/k)}{\log n}-\frac1M
 \ge u_0.
\end{equation}
Here $k/(n-2k+1)\le1$ for all sufficiently large $n$, and
$k>n^{3/5}$ makes $1/M\le1/1000$.  In particular
$\eta u<u^3$, so $\mathcal W(M,D)=Mu^3$ throughout this range.

We can therefore apply \eqref{rec:tuple-grid-second}.  Write
$\ell=\log n$, $b_0=\log(n/k)$, $\Delta=M-M'\ge0$, and
$\tau=k-(k_A+k_B)/2$.  If
$\varepsilon=(\log b-\log a)/2\in[0,(\log2)/2]$, then
\[
 \Delta=\tau\ell-\varepsilon(k_A-k_B),\qquad
 |\tau\ell-\Delta|\le k\log2.
\]
By \eqref{grid:perspective-middle-u}, $b_0/\ell\le u+1/M$.
Consequently
\begin{equation}\label{grid:perspective-middle-mass-loss}
 \frac{b_0\tau}{2}
 \le\frac{u\Delta}{2}+\frac12+\frac{k\log2}{2}
 \le\frac{u\Delta}{2}+2k.
\end{equation}
The identity $\Tr(RS)=\|SR\|_{S^2}^2$, the second-moment
overlap estimate, and rank allocation now give
\begin{equation}\label{grid:perspective-middle-indexed}
 \log s_j(Z)\le B-\min_{(R,S)}\left[
 c_0F_c+\frac12\{D-D_c-\tfrac12u\Delta-E_2\}_+\right],
 \qquad E_2=C_2k\ell^{3/4}+2k+B,
\end{equation}
where $C_2$ is the absolute constant in
\eqref{rec:tuple-grid-second}.

For clarity, the required tangent inequality follows directly
from $x^3\ge u^3+3u^2(x-u)$ for $x\ge0$.
Since $f(x)\ge x^3$ and $f(u)=u^3$, multiplication by each
$M_i$ and summation give
\begin{equation}\label{grid:perspective-middle-tangent}
 \begin{split}
 F_c&\ge3u^2D_c-2u^3M'
      =Mu^3-3u^2G,\\
 G&=D-D_c-\tfrac23u\Delta.
 \end{split}
\end{equation}
Zero-mass terms contribute zero.  Put
$H=D-D_c-\tfrac12u\Delta$.  The exact relation
$G=H-u\Delta/6\le H$ is favorable in the direction needed here.
If $H\le E_2$, the tangent inequality gives
$F_c\ge Mu^3-3u^2E_2$.  If $H>E_2$, it gives
\[
 c_0F_c+\tfrac12(H-E_2)
 \ge c_0Mu^3-3c_0u^2E_2
       +(\tfrac12-3c_0u^2)(H-E_2)
 \ge c_0Mu^3-3c_0u^2E_2,
\]
because $u\le2$ and $12c_0<1/4$.  Thus in either case
\eqref{grid:perspective-middle-indexed} implies
\begin{equation}\label{grid:perspective-middle-error}
 \log s_j(Z)\le-c_0\mathcal W(M,D)+B+12c_0E_2.
\end{equation}
The count \eqref{grid:perspective-profile-count} and
$k>n^{3/5}$ give the explicit bounds
\[
 \frac{B}{M}
 \le C_B\left(n^{-1/10}+\frac{n^{-1/20}}{\log n}\right),
 \qquad
 \frac{E_2}{M}
 \le C_2(\log n)^{-1/4}+\frac2{\log n}+\frac BM.
\]
Since $\mathcal W(M,D)=Mu^3\ge u_0^3M$, an absolute $K_2$
therefore satisfies
\begin{equation}\label{grid:perspective-middle-relative-error}
 \frac{B+12c_0E_2}{c_0\mathcal W(M,D)}
 \le K_2\left((\log n)^{-1/4}+n^{-1/20}\right).
\end{equation}
Again the bound is independent of $c$.

We can now make all choices in the induction.  Let $A\ge1$ be
an absolute constant and enlarge $d_0\ge5$ so that, whenever
$d>d_0$ and $n=2^d$, both right sides of
\eqref{grid:perspective-dense-relative-error} and
\eqref{grid:perspective-middle-relative-error} are at most
$Ad^{-1/8}$.  This is possible because their ratios to
$d^{-1/8}$ tend to zero.  Enlarge $d_0$ further to include all
previous absolute cutoffs, to have $2Ad_0^{-1/8}\le1/2$, and
to ensure
\begin{equation}\label{grid:perspective-top-index-error}
 12\le Ad^{-1/8}\mathcal W(M,\log(D_\lambda j))
\end{equation}
in both remaining regimes.  The last choice is uniform in $j$:
in the dense regime the weight is at least $\eta n^{99/100}$,
and in the middle regime it is at least
$u_0^3n^{3/5}\log n$.  The two error inequalities prove,
simultaneously for every $1\le i\le D_\lambda$,
\begin{equation}\label{grid:perspective-powered-final}
 s_i(Z)\le
 \exp\{-c_0(1-Ad^{-1/8})
              \mathcal W(M,\log(D_\lambda i))\}.
\end{equation}

Apply Lemma~\ref{grid:perspective-power-interpolation} and use
that a decreasing singular-value sequence satisfies
$s_j(CA)^j\le\prod_{i=1}^j s_i(CA)$.  We obtain
\[
 \log s_j(T_n|_{V_\lambda})
 \le-c(1-Ad^{-1/8})\frac1j
          \sum_{i=1}^j\mathcal W(M,\log(D_\lambda i)).
\]
No individual-singular-value interpolation is being assumed.
The same $12$-Lipschitz bound used above gives
\[
 \mathcal W(M,\log(D_\lambda i))
 \ge\mathcal W(M,\log(D_\lambda j))-12\log(j/i).
\]
Since $\log(j!)\ge j\log j-j$, the average of
$\log(j/i)$ over $1\le i\le j$ is at most $1$.
Thus the average weight is at least
$\mathcal W(M,\log(D_\lambda j))-12$.  In view of
\eqref{grid:perspective-top-index-error}, this proves the
inductive conclusion with exponent
\begin{equation}\label{grid:perspective-exponent-recursion}
 c_d=(1-2Ad^{-1/8})
       \min\{c_{\lfloor d/2\rfloor},c_{\lceil d/2\rceil}\}.
\end{equation}
The small-level estimate, considered at the start, also holds
with this exponent.

For $1\le d\le d_0$ choose a common $c_*\in(0,c_0]$ so small
that the desired inequality holds for every nontrivial type and
every singular index.  Such a choice exists by
Lemma~\ref{lem:finitegap}, because there are only finitely many
matrices, $\mathcal W\le8k\log n$, and zero operators need no
restriction on $c_*$.  Define $c_d=c_*$ on this finite range and
use \eqref{grid:perspective-exponent-recursion} thereafter.
This does not change the previously chosen cutoff.

Every multiplier in \eqref{grid:perspective-exponent-recursion}
is at least $1/2$ by the choice of $d_0$. The minimum selects a
branch of rounded halvings. Lemma~\ref{lem:dyadic-exponents},
with $\gamma=1/8$, therefore gives $\inf_d c_d>0$.
Taking this infimum as the theorem's exponent completes the proof.
\end{proof}
